% !TeX encoding = UTF-8
\documentclass[variant=docente,cover=institutional,mode=screen]{ua-engineering-v6-article}
% !TeX encoding = UTF-8
\UASetUniversity{Universidad Austral}
\UASetFaculty{Facultad de Ingeniería}
\UASetDepartment{Departamento de Computación}
\UASetCampus{Pilar}
\UASetCareer{Ingeniería Informática}
\UASetCommission{Comisión A}
\UASetProfessor{Equipo docente}
\UASetSemester{Segundo cuatrimestre 2026}
\UASetCourse{Sistemas Distribuidos}
\UASetDate{2026}


\UASetClassNumber{11}
\UASetClassTitle{Resiliencia distribuida}
\UASetDocKind{Guía docente}

\title{Clase 11 --- Guía docente}
\author{\UAProfessor}
\date{\UADate}

\begin{document}
\maketitle

\section{Objetivo pedagógico profundo}

La meta de esta clase es que el estudiante deje de pensar la resiliencia como una lista de patrones y la entienda como una disciplina de control de propagación de falla.

La idea que debe quedar instalada es:

\begin{center}
\textbf{una arquitectura resiliente no niega la falla; impide que la falla se propague sin control}
\end{center}

\section{Resultados de aprendizaje esperados}

Al finalizar, el estudiante debería poder:

\begin{itemize}
\item distinguir fallas transitorias de saturación persistente;
\item justificar cuándo retry ayuda y cuándo destruye;
\item explicar circuit breakers, backpressure y bulkheads como mecanismos complementarios;
\item diseñar una estrategia razonable de graceful degradation.
\end{itemize}

\section{Estructura sugerida para 90 minutos}

\subsection*{Bloque 1 (15 min): presentar la falla amplificada}
Abrir con un escenario realista:
\begin{itemize}
\item DB empieza a responder lento;
\item API hace retries;
\item sube la concurrencia;
\item el borde del sistema colapsa.
\end{itemize}

\paragraph{Objetivo didáctico}
Mostrar que el comportamiento incorrecto del sistema ante la falla puede ser peor que la falla original.

\subsection*{Bloque 2 (20 min): retries}
Explicar retry como herramienta útil pero peligrosa.

Trabajar:
\begin{itemize}
\item cuándo sirve;
\item cuándo no;
\item why exponential backoff;
\item why jitter.
\end{itemize}

\paragraph{Objetivo didáctico}
Corregir la intuición de que “más retry = más robustez”.

\subsection*{Bloque 3 (20 min): circuit breakers}
Explicar estados y propósito. Insistir en que el breaker no “cura” la dependencia: protege al sistema llamador y acelera feedback.

\paragraph{Objetivo didáctico}
Mostrar que el rechazo temprano puede ser una mejora sistémica.

\subsection*{Bloque 4 (15 min): backpressure y bulkheads}
Trabajar ambos juntos:
\begin{itemize}
\item backpressure controla entrada;
\item bulkheads aíslan recursos.
\end{itemize}

\paragraph{Objetivo didáctico}
Que el alumno vea la diferencia entre controlar presión y separar daño.

\subsection*{Bloque 5 (20 min): graceful degradation}
Discutir ejemplos de degradación aceptable e inaceptable según dominio.

\paragraph{Objetivo didáctico}
Mostrar que resiliencia no es solo infraestructura, sino también diseño funcional del producto.

\section{Qué explicar y por qué}

\subsection*{1. Retry}
Porque es el patrón más intuitivo y más frecuentemente mal aplicado.

\subsection*{2. Circuit Breaker}
Porque introduce explícitamente la idea de dejar de insistir.

\subsection*{3. Backpressure}
Porque sin control de demanda no hay estabilidad operativa sostenible.

\subsection*{4. Bulkheads}
Porque la separación de recursos evita contaminación cruzada entre flujos.

\subsection*{5. Graceful Degradation}
Porque hace visible que resiliencia también es semántica de servicio y no solo mecanismo técnico.

\section{Preguntas disparadoras recomendadas}

\begin{itemize}
\item ¿Qué pasa si todos reintentan al mismo tiempo?
\item ¿Qué recurso local se está consumiendo mientras la dependencia tarda?
\item ¿Qué debería dejar de funcionar primero y qué debería preservarse?
\item ¿Conviene rechazar rápido o esperar más?
\item ¿Qué parte del sistema no debería compartir pool con el resto?
\end{itemize}

\section{Errores frecuentes del alumnado}

\begin{itemize}
\item pensar que retry siempre ayuda;
\item usar circuit breaker sin fallback semántico;
\item hablar de backpressure solo como “colas grandes”;
\item creer que bulkheads son solo un detalle de infraestructura;
\item confundir graceful degradation con respuestas incorrectas o engañosas.
\end{itemize}

\section{Ejercicio en vivo sugerido}

Tomar un flujo de checkout con:
\begin{itemize}
\item API;
\item Auth;
\item Payments;
\item Recommendations.
\end{itemize}

Plantear que Payments empieza a degradarse.

Preguntar al curso:
\begin{itemize}
\item dónde poner timeouts;
\item si conviene retry;
\item cuándo abrir breaker;
\item qué recursos aislar;
\item qué funcionalidad degradar.
\end{itemize}

\paragraph{Objetivo}
Que el alumno piense resiliencia como arquitectura de contención y no como una sola librería.

\section{Momento crítico de la clase}

El punto más importante aparece cuando el estudiante comprende que:

\begin{center}
\textbf{la respuesta correcta a una dependencia degradada no siempre es insistir más, sino insistir menos y aislar mejor}
\end{center}

Ese punto debe subrayarse, porque rompe una intuición muy frecuente y muy dañina.

\section{Criterio de éxito docente}

La clase salió bien si el estudiante puede:

\begin{itemize}
\item explicar por qué retry puede empeorar saturación;
\item describir correctamente estados de un breaker;
\item distinguir backpressure de bulkheads;
\item proponer una degradación razonable para un dominio concreto;
\item diseñar una estrategia integrada de resiliencia.
\end{itemize}

\section{Conexión con la clase siguiente}

Esta clase deja preparada una pregunta natural:

\begin{quote}
Si ya sabemos diseñar, observar y volver resiliente el sistema, ¿cómo validamos experimentalmente que realmente se comporta bien bajo fallas?
\end{quote}

Eso conecta naturalmente con:
\begin{itemize}
\item chaos engineering;
\item fault injection;
\item testing distribuido;
\item validación experimental de resiliencia.
\end{itemize}

\begin{uakeybox}
Pregunta puente: ¿cómo pasamos de confiar en nuestros mecanismos de resiliencia a demostrar empíricamente que funcionan?
\end{uakeybox}


\section{Plan de conducción ampliado}
La clase debe alternar explicación, predicción y contraste. Antes de revelar cada mecanismo, pedir al grupo que anticipe qué puede salir mal; después, volver al mismo escenario y preguntar qué propiedad mejora y qué costo aparece. Cada bloque conceptual debe cerrar con una decisión de diseño observable.
\subsection*{Secuencia recomendada}
\begin{enumerate}
\item Recuperar el prerequisito de la clase anterior.
\item Presentar un caso mínimo que falle y solicitar hipótesis.
\item Formalizar el concepto central de Resiliencia distribuida, distinguiendo seguridad, progreso y supuestos.
\item Resolver un ejemplo completo en pizarra, incluyendo una falla intermedia.
\item Comparar una alternativa de diseño y explicitar trade-offs.
\item Cerrar con una pregunta del Trabajo Final y una pregunta de defensa oral.
\end{enumerate}
\section{Diagnóstico de comprensión durante la clase}
No preguntar solamente ``¿se entendió?''. Usar preguntas discriminantes: ``¿qué cambia si el mensaje llega dos veces?'', ``¿qué propiedad sigue siendo cierta si cae un nodo?'', ``¿esto demuestra seguridad o progreso?'', ``¿qué observa un cliente externo?''. Si el estudiante responde solo con nombres de patrones, pedir una ejecución concreta.
\section{Criterios para corregir ejercicios}
Una respuesta excelente declara supuestos, construye una secuencia verificable, nombra la garantía con precisión, reconoce un costo y conecta la decisión con el dominio. Penalizar afirmaciones absolutas sin condiciones.
\section{Vinculación con el Trabajo Final Integrador}
Reservar entre 5 y 10 minutos para que cada grupo registre una decisión con: problema, alternativa elegida, alternativa descartada y razón. Esto crea trazabilidad entre las clases y las entregas acumulativas.
\section{Lectura docente previa}
Revisar la bibliografía indicada en los apuntes y preparar al menos un contraejemplo que no aparezca literalmente en las diapositivas. El objetivo es poder responder preguntas de frontera sin convertir la clase en una lectura del deck.

\section{Hito del Trabajo Final: Etapa 2}
Al inicio de la clase corresponde la entrega acumulativa de la Etapa 2. La devolución debe revisar la coherencia entre el punto de consenso detectado antes, la solución concreta elegida, el modelo de consistencia, la estrategia de datos, las transacciones y el modelo de fallas. El prototipo debe demostrar un mecanismo real, no una captura estática.

\section{Arquitectura didáctica v9 para 180 minutos}
\subsection*{Bloque 1 - desafío inicial (20 min)}
Presentar una situación donde aparezca dependencia lenta, retry storm, cola sin límite, pool agotado y cascada. Pedir predicciones individuales antes de introducir vocabulario. El objetivo es hacer visibles los supuestos intuitivos y recuperar el prerequisito de la clase anterior.

\subsection*{Bloque 2 - construcción conceptual (35 min)}
Formalizar retry budgets, backoff, jitter, circuit breakers, backpressure, bulkheads y graceful degradation. Separar seguridad, progreso y experiencia del usuario. Explicar no solo \emph{qué} hace cada mecanismo, sino \emph{por qué} existe y qué supuesto permite que funcione.

\subsection*{Bloque 3 - modelo y figura (35 min)}
Construir en pizarra una línea temporal, grafo, log, máquina de estados o tabla de decisiones. Recorrer una ejecución nominal y una adversarial. La representación debe quedar como referencia para los apuntes y el ejercicio principal.

\subsection*{Pausa (10 min)}
Recoger una pregunta anónima o un punto de confusión. Elegir una pregunta que permita distinguir síntoma, causa y propiedad.

\subsection*{Bloque 4 - caso trabajado con falla (40 min)}
Aplicar el mecanismo al caso conductor. Introducir una falla a mitad de ejecución y detenerse en cada transición: quién sabe qué, qué se persiste, qué garantía sigue vigente y qué observa el cliente.

\subsection*{Bloque 5 - taller de decisión (25 min)}
Los grupos aplican P-F-G-S-T a su Trabajo Final. Deben producir un ADR breve con alternativa elegida, alternativa descartada y evidencia pendiente. La evidencia de esta clase es curvas de carga, colas, breaker state y degradación del camino crítico.

\subsection*{Bloque 6 - cierre y exit ticket (15 min)}
Cada estudiante responde: propiedad principal, supuesto decisivo, costo aceptado y condición bajo la cual cambiaría de diseño. Usar las respuestas para iniciar la clase siguiente.

\section{Qué explicar, por qué y cómo verificarlo}
\begin{itemize}
\item Explicar el problema antes del producto, porque reconocer la familia de problema es el resultado cognitivo central.
\item Explicar el invariante, porque permite evaluar configuraciones y no solo repetir un flujo nominal.
\item Explicar el límite, porque dependencia lenta, retry storm, cola sin límite, pool agotado y cascada puede invalidar supuestos o reducir progreso.
\item Verificar comprensión mediante curvas de carga, colas, breaker state y degradación del camino crítico, no mediante una definición aislada.
\end{itemize}

\section{Preguntas socráticas}
\begin{itemize}
\item ¿Qué sabe cada actor en este punto de la ejecución?
\item ¿Qué propiedad sigue siendo cierta si la respuesta nunca llega?
\item ¿Qué parte es safety y cuál es liveness?
\item ¿Qué alternativa más simple resolvería el problema si relajamos una garantía?
\item ¿Qué observación refutaría nuestra hipótesis de diseño?
\end{itemize}

\section{Errores previsibles y estrategia de corrección}
\begin{itemize}
\item Si el alumno responde con un nombre de tecnología, pedir actores, mensajes y estados.
\item Si confunde timeout con falla, construir dos ejecuciones indistinguibles.
\item Si promete todas las garantías, pedir comportamiento bajo partición o saturación.
\item Si mide solo rendimiento, preguntar cómo evidencia la propiedad semántica.
\item Si la solución es excesiva, pedir la alternativa mínima y el costo marginal de la garantía fuerte.
\end{itemize}

\section{Criterios de evaluación formativa}
Una respuesta excelente declara supuestos, recorre una ejecución, nombra con precisión la garantía, reconoce el costo y propone evidencia reproducible. Una respuesta suficiente identifica el mecanismo y el trade-off principal. Una respuesta insuficiente usa slogans, omite fallas o confunde producto con propiedad.

\section{Preparación docente y bibliografía}
Lectura previa recomendada: Nygard, Release It!; Google SRE Workbook; patrones de backpressure y load shedding. Preparar un contraejemplo adicional y una variante del caso en la que cambie el costo del error; esto evita convertir la clase en una lectura de diapositivas.

\end{document}
